% 第 27 章　气体是怎样顶住容器的（完整章）
\chapter{气体是怎样顶住容器的}

\step{反常识开场}

\begin{mainline}
找一支干净的注射器（去掉针头），把活塞推到底，然后用手指死死堵住针尖的小孔。现在把活塞往回拉——你会感到一股明显的“吸力”在跟你抢活塞，拉到一半松手，活塞“嗖”地一下被吸回去大半截。换一种玩法：把活塞先拉到中间，堵住小孔，然后往里推。越推越费力，像里面藏了一根弹簧；松手，活塞又被弹了回来。

停下来想一想这件事有多奇怪。针筒里有什么？什么都没有——没有弹簧，没有磁铁，没有橡皮筋，只有空气。而“空气”在我们的日常语言里几乎就是“什么都没有”的代名词：空着手、空房间、一场空。可就是这“什么都没有”，既能像弹簧一样把活塞顶回来，又能像手一样把活塞拽回去。它凭什么？

更奇怪的还在后面。把这只堵住口的针筒整个泡进热水里，别碰活塞，你会看到活塞自己缓缓向外移动——热水没有碰活塞，它只是在加热“里面的空气”，而空气就把活塞顶走了。热是怎么变成推力的？

类似的事情其实每天都在发生。从平原开车上高原，后备箱里密封的薯片袋会鼓成一只小枕头，鼓鼓囊囊仿佛随时要炸开——袋子是密封的，里面的空气一克都没有增加，它为什么在山上突然变得“气鼓鼓”？给自行车打气，打到后面每一下都像在跟一堵墙较劲，可轮胎的体积几乎没怎么变大，打进去的空气到底“挤”在哪里？

空气看不见、摸不着、几乎感觉不到重量，却能顶活塞、鼓袋子、硬抗打气筒。这一章的任务，就是给这个隐形的巨人画一幅像：它到底是谁，它用什么办法顶住容器，以及——为什么偏偏是“加热”能让它变得更有劲。
\end{mainline}

\step{先猜一猜}

照例，先下注再开牌。请把下面三题的答案写下来，越具体越好。

第一题。堵住口的针筒，活塞拉到一半松手，活塞被“吸”回去。现在想象一个理想化的版本：针筒里不是空气，而是真正的真空——一个分子都没有，外面仍然是普通空气。松手后活塞会怎样？（a）停在原地不动，因为里面没有东西推它；（b）被一路推到底，而且推得比有空气时更狠；（c）先往里走一点再停住。

第二题。两只完全相同的气球，充到一样大，一只装空气，一只装氦气（就是能让声音变尖的那种轻气体），温度相同。哪只气球内部的压强大？（a）装空气的大，因为空气分子更重，撞得更疼；（b）装氦气的大，因为氦气分子轻、跑得快，撞得更勤；（c）一样大。

第三题。一只密封的刚性玻璃瓶（不会变形），里面就是普通空气。把它从冰箱（约 4 摄氏度）拿到室温（约 25 摄氏度）的桌上放着。瓶内空气的压强会怎样？（a）不变，因为瓶子是密封的，空气跑不掉；（b）变大；（c）变小。如果你选（b），再追一句：变大的比例大概是多少——百分之一量级，还是百分之十量级？

写下你的答案和理由。本章结尾我们会回来核对。猜错不丢人——这一章要修理的，正是“看不见的东西做不了大事”这条直觉。

\step{动手实验}

\begin{experiment}[针筒里的隐形弹簧]
\textbf{材料}：一支 10 或 20 毫升的塑料注射器（药房有售，务必不使用针头）、一小盆热水、一小盆冷水（或冰块）。

\textbf{步骤一（压）}：把活塞拉到大约一半刻度，用手指紧紧堵住针尖小孔，另一只手往里推活塞。感受阻力如何变化：推得越深入，阻力越大，而且不是均匀地变大——压到一半体积时，你会明显感觉“再压一倍”几乎做不到。松手，活塞弹回，但不是回到原位，而是停在一个比原位稍浅的位置（手指堵住时漏不掉的那部分空气被压缩过，再平衡时稍有差别）。

\textbf{步骤二（拉）}：重新拉到中间，堵住孔，往外拉活塞。同样越拉越费力，像里面有一只手在往回拽。松手，活塞被“吸”回去。

\textbf{步骤三（加热）}：活塞推到约三分之一处，堵孔，把针筒水平固定好（比如用橡皮筋绑在筷子上架着），把热水浇在针筒外壁上，或者在热水盆里浸一分钟。盯住活塞：它会自己向外走。换冷水浇上去，活塞又慢慢缩回来。整个过程中你的手没有碰活塞。

\textbf{步骤四（对照）}：把针筒吸满水，堵住孔，再用力推。水的阻力大得多，而且活塞几乎纹丝不动——水不像空气那样能被压缩出明显的一截。同样是堵住孔往里推，空气像弹簧，水像石头。记住这个差别。

\textbf{记录}：至少写下三个观察：压缩和拉伸时阻力如何随位置变化；加热和冷却时活塞往哪个方向走；水和空气的表现差在哪里。暂时不需要解释，本章后面会一条一条拆开。
\end{experiment}

还有一个零成本观察：找一个塑料吸盘挂钩，按在光滑的瓷砖上，挂上一串钥匙。多数人会这样描述它：“吸盘吸住了墙。”请在心里给这句话画个问号——“吸”是谁在吸？朝哪个方向用力？这个问号我们先存起来，它会成为本章最好的反面教材。

\step{旧模型遇到麻烦}

在动手拆直觉之前，先公平地陈述一下旧模型。面对针筒实验，一个完全合理的解释是：空气是一种连续的、有弹性的“果冻”，压它就反抗，拉它就回缩，像一块透明橡皮。这个模型听起来没毛病，而且历史上人们确实这样想了很久——把空气当成一种连续介质，用“弹性”二字打发掉所有问题。

但麻烦一个接一个地来了。

麻烦一：弹性果冻解释不了“加热”。橡皮筋靠分子间的拉力储存弹性，可你把针筒泡进热水，活塞自己往外走——果冻为什么会因为变热而膨胀得更想顶人？连续模型只能说“它就是这样”，给不出任何机制。更糟的是麻烦一的定量版本：实验早就发现，一定量的空气，压强乘体积在温度不变时几乎是个常数（压成一半体积，压强正好翻倍），温度升高一度，压强又按一个干净的比例上涨。果冻模型说不出这些数字为什么干净。

麻烦二：四面八方同时存在的压强。把一只玻璃杯装满水，盖上硬纸片，倒过来，纸片不掉，水不洒——空气从\emph{下方}把纸片和水托住了。一个开口朝下的杯子，里面的水面照样被空气顶着。如果空气只是“被压才会反抗”的果冻，它凭什么主动向上托？果冻模型里压强是“你挤它它才挤你”，可覆杯实验里没有任何人去挤空气，是空气自己在每一平方米面积上托着约十牛每平方厘米的东西，方向任选：向上、向下、向左、向右，全都一样大。

麻烦三：也是最致命的一条——“空”针筒里的弹簧。如果空气是连续果冻，那把空气抽走之后，果冻没了，针筒里应该什么都不会发生。可事实上（这正是第一题的场景），真空的针筒里活塞会被外部空气一路推到底。真正奇怪的事情翻转了：需要解释的不是“真空怎么吸活塞”，而是“空气怎么推活塞”——果冻模型连谁在给谁施力都说不清楚，它只会说“有吸力”，而“吸力”这个词把方向搞反了：吸盘不是被墙“吸住”的，是被墙外的空气“压住”的；针筒活塞不是被真空“吸回去”的，是被外面的空气“推回去”的。

三条麻烦的共同诊断：把空气当成一整块连续的东西，就永远回答不了“它凭什么施力”“为什么加热会变强”“数字为什么那么干净”。连续模型不是全错——后面我们会看到，在大多数工程计算里把空气当连续介质完全够用——但它丢掉了空气的内部结构，也就丢掉了“为什么”。要找回“为什么”，得换一个完全不同的画面：空气不是一块果冻，而是一大群东西。

\step{发明新概念}

\begin{mainline}
\textbf{新概念一：气体是一大群永不停歇的微小粒子。}这是十九世纪逐渐成形的分子动理论，画面极其简单：空气由数量大到荒谬的微小分子组成，每个分子都很轻、很小，彼此之间的平均距离远大于自身尺寸，它们在容器里沿着直线乱飞，撞墙、反弹、偶尔互相撞上，再接着乱飞，永不停歇。注意这个画面里没有“弹性物质”，没有任何连续的果冻——只有粒子、运动和碰撞。

\textbf{新概念二：压强是无数碰撞的集体平均效果。}单个分子撞一下墙，就像一粒小雨点砸在伞上，力极小、时间极短，谁都感觉不到。可是当每秒钟有难以想象的那么多次撞击均匀地落在同一面墙上时，无数次微小而短促的推，平均成一股稳定、持续、光滑的压力。它像什么？像暴雨打在伞面上：单个雨滴是“啪、啪、啪”的一下下撞击，雨密到一定程度，你感觉到的只是一股均匀往下压的力。它不像什么？它不像果冻的弹性——果冻靠内部分子被拉伸后的回缩力反抗，而气体的压力靠的是\emph{撞击}，粒子撞上来、弹回去，把动量交给墙。这个区别不是修辞，它马上会解释旧模型回答不了的一切。

先看它如何一举拆掉三条麻烦。压成一半体积，同样多的分子挤在一半的空间里，撞墙的频率正好翻倍——压强正好翻倍，那个干净的数字有了来源。加热让分子跑得更快，每次撞墙撞得更狠、撞得更勤——压强随温度上涨，针筒泡在热水里活塞被顶走，有了机制。而每个分子撞墙的方向是四面八方乱来的，平均下来朝任何方向的碰撞都一样多——覆杯实验里空气向上托纸片，和向下压桌面，用的是同一批分子的同一种乱撞。至于真空针筒：里面没有分子，里面那一侧没有任何碰撞；外面的空气分子照样每秒亿万次撞活塞，于是活塞被外面推着走，直到撞到底。根本没有什么“吸力”，从来都是\emph{两边的碰撞之差}在推活塞。

\textbf{新概念三：温度是分子运动剧烈程度的度量。}上一章我们从热平衡出发引入了温度，但还没有说温度“是什么”。现在分子动理论给出答案：对气体而言，温度度量的正是分子无规则运动的平均动能——温度越高，分子的平均动能越大。这句话要立刻配两条警告。第一，温度是\emph{一大群}分子的统计性质，说“这一个分子的温度是多少”没有意义，就像说“某一位观众的平均掌声”没有意义。第二，是\emph{平均}动能：同一温度下，分子有快有慢，撞来撞去不断交换速度，温度只盯住那个平均值。
\end{mainline}

这套画面漂亮得可疑。十九世纪的怀疑者反问得很合理：分子看不见摸不着，你们这套“粒子乱撞”说得好听，证据呢？是不是又一个讲得通的神话？回答这个问题的，是科学史上最漂亮的一次“假说变实证”——布朗运动。

1827 年，植物学家布朗在显微镜下观察悬浮在水中的花粉颗粒，看到它们一刻不停地做着无规则的抖动，方向随机、永不停歇。这是\textbf{实验事实}：任何人用一台显微镜和一点花粉（或者稀释的墨汁、奶粉）都能复现，颗粒越小抖得越凶，水越热抖得越凶。接下来的部分是\textbf{解释}：1905 年，爱因斯坦提出，花粉颗粒之所以抖动，是因为它被四面八方的水分子不停撞击——颗粒足够小时，某一瞬间左边撞来的分子恰好比右边多几个，颗粒就被踹了一脚；下一瞬间方向又随机换了。注意，花粉颗粒本身比分子大几十万倍，单个分子的撞击根本挪不动它，能被看见，全靠涨落——碰撞次数的微小不平衡。再往后是\textbf{检验}：佩兰用显微镜逐帧追踪了大量颗粒的轨迹，发现抖动的定量规律和爱因斯坦基于分子假说算出的预言严丝合缝，还顺手测出了分子的数目尺度。到这一步，“原子和分子真实存在”从一种好用的假说，变成了可检验、已检验的对象。要特别强调一句：布朗运动本身不是“分子的运动”，你看到的是花粉的抖动；分子的存在和撞击是对这抖动的解释——只不过这个解释做出了精确的定量预言，而且全部兑现，好到让对手认输。

\step{一句话模型}

\begin{oneline}
气体不是连续的东西，而是一大群乱飞的分子；压强就是无数分子撞击器壁的平均效果——单位面积上每秒接收到的碰撞越频繁、每次撞得越狠，压强就越大；温度越高，分子平均动能越大，撞得也就越狠。容器之所以被“顶住”，是因为墙内外两侧的碰撞不一样多。
\end{oneline}

用这一句话检查前面的每个现象：针筒被压缩后内部碰撞变密，活塞被顶回来；真空针筒内部零碰撞，外部碰撞把活塞推到底；热水让分子加速，活塞被顶走；薯片袋上了高原，外部空气变稀薄、外部碰撞变少，内部碰撞占了上风，袋子鼓起来。一句话，全罩住了。

\step{数学翻译器}

直觉已经到位，现在把它翻译成能算数的语言。分两步走：先算“一个分子撞一次墙交出多少东西”，再算“亿万次碰撞平均成多大的压强”。

\textbf{第一步：一次碰撞。}设一个分子质量为 \(m\)，以速度 \(v\) 垂直撞上一面墙，原速率弹回。它的速度从 \(+v\) 变成 \(-v\)，动量变化了 \(2mv\)。根据动量守恒（第 10 章的内容），墙收到了 \(2mv\) 的动量。碰撞越频繁（单位时间撞上墙的分子越多）、每次交出的动量越大（\(m\) 大、\(v\) 大），墙受到的平均力就越大。

\textbf{第二步：平均成压强。}压强 \(p\) 定义为单位面积上的力：\(p = F/S\)。把盒子里 \(N\) 个分子的全部碰撞加起来——考虑到分子速度有各个方向、快慢不一，要用平均值——最后得到的是一条结构非常漂亮的比例关系：
\[
p \;=\; \frac{N}{V}\cdot \frac{1}{3}\, m \overline{v^{2}},
\]
其中 \(V\) 是容器体积，\(\overline{v^{2}}\) 是分子速率平方的平均值，那个 \(1/3\) 来自“速度平均分摊到前后、左右、上下三个方向”。先别管推导细节（下面数学桥层里补全），先读这个式子在说什么：压强正比于\emph{分子的密集程度}（\(N/V\)）乘以\emph{单个分子的动能}（\(m\overline{v^2}/2\) 的两倍）。这正是上一节一句话模型的符号版：撞得越密、撞得越狠，压强越大。

\begin{center}
\begin{tikzpicture}[scale=1.05]
  % 容器
  \draw[thick] (0,0) rectangle (5.4,2.8);
  % 活塞
  \fill[gray!25] (3.8,0) rectangle (4.2,2.8);
  \draw[thick] (3.8,0) rectangle (4.2,2.8);
  \node[above] at (4.0,2.8) {活塞};
  % 分子与速度箭头
  \fill (0.6,0.5) circle (1.5pt); \draw[-{Stealth}] (0.6,0.5) -- (1.2,0.9);
  \fill (1.5,2.1) circle (1.5pt); \draw[-{Stealth}] (1.5,2.1) -- (2.2,1.8);
  \fill (2.6,0.8) circle (1.5pt); \draw[-{Stealth}] (2.6,0.8) -- (2.0,1.3);
  \fill (0.9,1.6) circle (1.5pt); \draw[-{Stealth}] (0.9,1.6) -- (1.6,1.5);
  \fill (2.2,2.3) circle (1.5pt); \draw[-{Stealth}] (2.2,2.3) -- (1.7,2.0);
  \fill (3.1,1.7) circle (1.5pt); \draw[-{Stealth}] (3.1,1.7) -- (3.5,1.2);
  \fill (1.9,0.4) circle (1.5pt); \draw[-{Stealth}] (1.9,0.4) -- (2.5,0.7);
  % 内部碰撞把活塞向右推
  \draw[-{Stealth},very thick,blue!60!black] (3.1,0.35) -- (3.75,0.35);
  \node[below,blue!60!black] at (3.45,0.3) {\small 内部碰撞};
  % 外部大气把活塞向左推
  \draw[-{Stealth},very thick,red!60!black] (5.6,1.4) -- (4.25,1.4);
  \node[right,red!60!black] at (5.6,1.4) {\small 外部大气的碰撞};
\end{tikzpicture}
\end{center}

活塞两侧永远在挨撞：内侧是气体分子，外侧是大气分子。活塞往哪边走，不取决于任何一边“有多想推”，而取决于两侧的碰撞之差。下面整个数学翻译，就是把这幅图变成能算数的式子。

\textbf{立刻用特殊情况检查它。}体积不变、分子数 \(N\) 翻倍——压强翻倍：给轮胎多打气，胎内气压确实上涨，对。\(N\) 减到零（真空）——压强为零：真空针筒内部不推活塞，对。所有分子停下来（\(\overline{v^2}=0\)）——压强为零：一群静止的粒子堆在容器底部，确实顶不住任何东西，对。体积翻倍而分子数不变——压强减半：这正是“压成一半体积压强翻倍”那条古老的实验定律（玻意耳定律），我们的式子自动包含了它。

\textbf{第三步：把温度请进来。}上一节说温度度量分子的平均动能，现在把这句话写成等式：对气体分子在容器中的无规则运动，
\[
\frac{1}{2} m \overline{v^{2}} \;=\; \frac{3}{2}\, k T,
\]
其中 \(T\) 是用绝对温标（开尔文）表示的温度，\(k\) 是一个新的自然常数——玻尔兹曼常数，
\[
k \approx 1.38\times 10^{-23}\ \text{J/K}.
\]
这个常数值得停下来看一秒钟。等式左边是\emph{单个分子}的动能，微观世界的量；右边是\emph{温度}，温度计读出来的宏观世界的量。\(k\) 就是这两个世界之间的汇率：每升高一开尔文，每个分子的平均动能增加 \(1.5\times 1.38\times 10^{-23}\) 焦耳。数字小得可怜，因为单个分子小得可怜——这正是我们感觉不到单个分子的原因，也是我们需要万亿亿个分子才能感到一股稳定压力的原因。

把两步合起来，就得到了本章的中心公式——\textbf{理想气体状态方程}：
\[
pV \;=\; NkT.
\]
读法：一定量（\(N\) 固定）的气体，压强乘体积正比于绝对温度。它也常被写成 \(pV = nRT\) 的形式，那是把分子按“摩尔”打包计数的等价写法（每摩尔的分子数是个固定的大数，\(R = k\) 乘上这个大数）。

\textbf{再检查一遍状态方程。}\(T\) 不变时 \(pV\) 是常数——玻意耳定律，老朋友。\(V\) 固定时 \(p\) 正比于 \(T\)：密封玻璃瓶从 277 开尔文（约 4 摄氏度）拿到 298 开尔文（约 25 摄氏度），压强变成原来的 \(298/277 \approx 1.08\) 倍——第三题的答案是“变大，约百分之八”，注意比例要用开尔文温度算，用摄氏度会得到荒谬的六倍。\(p\) 固定时 \(V\) 正比于 \(T\)：冷却气球，气球变瘪，放冰箱冷藏室里的气球确实会缩小一圈。零温外推下去 \(p\) 趋向零——经典图像里分子全部停下，压强消失；真实世界在到达那一步之前气体早就液化甚至凝固了，这是模型的边界，下一节细说。

\textbf{一个带真实数字的估算：空气分子跑得有多快？}由 \(\frac{1}{2}m\overline{v^2}=\frac{3}{2}kT\)，室温（约 300 开尔文）下一个氮气分子（质量约 \(4.7\times10^{-26}\) 千克）的典型速率是
\[
\sqrt{\overline{v^2}} \approx \sqrt{\frac{3\times 1.38\times 10^{-23}\times 300}{4.7\times 10^{-26}}} \approx 5\times 10^{2}\ \text{m/s}.
\]
每秒五百米，比声音在空气中的传播速度还快一截。你此刻周围每一个空气分子都在以子弹般的速度乱飞，只是因为它们极其微小、撞击在你皮肤上的每一粒都轻若无物，你才毫无知觉。而它们数量惊人：一立方厘米的室内空气约有 \(2.5\times10^{19}\) 个分子。每秒钟撞在你手背一平方厘米皮肤上的分子数，粗估在 \(10^{23}\) 量级——万亿亿次。稳恒的气压，就是这么多次微小撞击平均出来的。

\textbf{再来一个估算，顺带给直觉上一课：大气压有多大？}地面附近的大气压约为 \(10^5\) 牛每平方米，也就是每平方厘米约十牛——相当于在你的指甲盖上放一袋一千克的水果。你的全身皮肤面积约 \(1.5\) 平方米，承受的大气压力合计约 \(1.5\times10^5\) 牛，相当于一头十几吨重的鲸鱼压在身上。为什么我们没被压扁？因为体内外的碰撞是平衡的：你体内的液体和气体从内侧以同样的压强向外撞，外侧大气从外侧向内撞，两边抵消。这和真空针筒的道理一模一样——起作用的永远是碰撞之差。也正因为如此，潜水员每下潜十米，外部压强增加约一个大气压，而体内压强必须通过呼吸加压气体同步跟上，否则胸腔会被这个“差”压伤：潜水病的全部安全教育，本质都是碰撞之差的管理学。

\begin{center}
\begin{tikzpicture}[scale=1.0]
  \draw[-{Stealth}] (0,0) -- (4.8,0) node[right] {$V$};
  \draw[-{Stealth}] (0,0) -- (0,3.2) node[above] {$p$};
  \draw[domain=0.85:4.4,smooth,thick] plot (\x,{2.4/\x});
  \fill (1.2,2.0) circle (2pt) node[above left] {\small 体积 \(V_0\)，压强 \(p_0\)};
  \fill (2.4,1.0) circle (2pt) node[above right] {\small 体积 \(2V_0\)，压强 \(p_0/2\)};
  \node at (3.6,2.4) {\small 温度不变：\(pV=\) 常数};
\end{tikzpicture}
\end{center}

\begin{mathbridge}
\textbf{从一次碰撞到 \(pV=\frac{1}{3}Nm\overline{v^2}\) 的完整推导。}取边长为 \(L\) 的立方盒子，先看一个分子在 \(x\) 方向的分运动：它以速率 \(v_x\) 在两堵相对的墙之间来回，撞一次墙动量改变 \(2mv_x\)，往返一次的路程是 \(2L\)，所以撞上同一面墙的时间间隔是 \(2L/v_x\)，平均每秒撞 \(v_x/(2L)\) 次。这面墙从这个分子身上收到的平均力，等于“每次的动量”乘“每秒的次数”：
\[
F_x = 2mv_x \cdot \frac{v_x}{2L} = \frac{mv_x^2}{L}.
\]
\(N\) 个分子总贡献 \(F = \frac{m}{L}\sum v_x^2 = \frac{mN}{L}\overline{v_x^2}\)。压强等于力除以墙的面积 \(L^2\)：\(p = \frac{mN}{L^3}\overline{v_x^2} = \frac{N}{V}m\overline{v_x^2}\)。最后一步用对称性：分子乱飞没有任何偏爱方向，\(\overline{v_x^2}=\overline{v_y^2}=\overline{v_z^2}\)，而 \(\overline{v^2}=\overline{v_x^2}+\overline{v_y^2}+\overline{v_z^2}=3\overline{v_x^2}\)，于是
\[
pV = \frac{1}{3}Nm\overline{v^2}.
\]
再代入 \(\frac{1}{2}m\overline{v^2}=\frac{3}{2}kT\)，即得 \(pV=NkT\)。整条推导只用了动量、平均和“三个方向对称”这三样东西。
\end{mathbridge}

\begin{challenge}
\textbf{为什么是“温度”而不是别的什么量进入状态方程？}更深一层的问法是：两个温度不同的气体混在一起，是什么保证它们最终达到同一个 \(T\)？答案藏在碰撞的统计规律里：分子间随机碰撞不断交换能量，会趋向一种最可几的速度分布（麦克斯韦分布），而这个分布里唯一的自由度恰好就是平均动能——也就是温度。这是第 31 章统计物理的入口：宏观状态量（温度、压强）之所以稳定而简单，是因为它们是天文数字次随机碰撞平均后的幸存者，涨落被数量级压掉了。配分函数、玻尔兹曼分布这些更锋利的工具，会在那里等着。
\end{challenge}

\step{模型的边界}

理想气体模型是一份合同，签字前要看清条款。它假设了两件事：分子自身的体积可以忽略（相对于容器的空旷程度）；分子之间除了碰撞瞬间外没有相互作用力（既不相吸也不相斥，各飞各的）。在这两条成立的地方，\(pV=NkT\) 精确得惊人——常温常压下的空气，误差不到百分之一。

条款在两种场合失效。第一种：\textbf{压得太狠}。把气体压缩到分子间距只有分子尺寸的十来倍以内，“分子本身有体积”这件事就赖不掉了——可压缩空间比标称体积小，实际压强比状态方程预言的更大。再往下压，气体开始液化：状态方程说压强可以无限涨，真实气体却用相变投了反对票。液化气钢瓶就是例子：瓶里大部分是液体，上方是蒸气，压强由温度而不是体积决定，\(pV=NkT\) 对瓶内整体根本不适用。第二种：\textbf{冷到分子间吸引力显形}。分子低速飞行时，彼此之间微弱的吸引力有机会把它们拉得靠近一点，压强会比理想值偏小；再冷下去，吸引力直接获胜——凝结。真实气体的修正版本（范德瓦尔斯方程）给体积和压强各加一项修正，正好对应这两条失效，那是物理化学的内容，这里只标记路牌。

还有两类典型的\textbf{误用}，比数值误差更危险。误用一：把温度套在单个分子上。“这个分子很热”是无意义的话——温度是平均出来的统计量，一个分子只有动能，没有温度。误用二：把压强想象成“气体被什么东西压着”。封闭气球里的压强不是谁施加给气体的，是气体分子自己撞出来的；内外碰撞之差才决定膜的受力。误用三（这条来自旧模型的阴魂）：说吸盘“吸”在墙上。真空中没有东西提供吸力，是墙外每平方厘米约十牛的大气压力把吸盘按在墙上；把吸盘连墙一起搬进真空室，它自己就掉了。

\begin{thought}[极端尺度：只剩一百个分子的容器]
把容器一路缩小、抽稀，直到里面只剩一百来个分子。墙上挂一只足够灵敏的压强计，读数会是什么样？不再是一条平稳的直线，而是一根上蹿下跳的毛刺曲线：某一瞬间恰好有六个分子撞在这面墙上，读数冲高；下一瞬间只来了两个，读数跌下去。“压强”开始闪烁。这暴露了一个平时被数量级藏起来的事实：压强本质上是平均值，它的稳定性来自碰撞次数的天文数字——每秒 \(10^{23}\) 次碰撞，相对涨落小到十万亿分之一量级，仪器永远测不到；只剩一百个分子时，涨落和平均值一样大，“压强”这个概念本身开始失去意义。宏观规律的稳定性是统计出来的，不是定律禁止它闪烁。这个思想实验是通往第 31 章的桥：在那里，涨落、分布和大数定律将成为主角。
\end{thought}

\step{回头解释开场现象}

现在回到针筒，把开场的每个谜团逐一销案。

\textbf{推活塞为什么像压弹簧？}堵住小孔后，针筒内的分子数 \(N\) 固定。活塞往里推，\(V\) 变小，由 \(p=NkT/V\)，内部压强按比例上升——压到一半体积，内部压强约两倍大气压。活塞内侧面承受的碰撞比外侧面（始终是一个大气压）多出一截，这个碰撞之差就是顶你手的力，压得越狠顶得越凶。松手后，内侧占上风的碰撞把活塞一路推回内外压强平衡的位置。所谓“隐形弹簧”，其实是每秒万亿亿次碰撞的差值。

\textbf{拉活塞为什么像被吸回去？}拉活塞使 \(V\) 变大，内部压强跌破大气压，这一次是\emph{外侧}碰撞占上风，把活塞推回去。从头到尾没有任何东西在“吸”——第一题的答案由此明确：真空针筒里内部零碰撞，外部大气压以每平方厘米约十牛的力气把活塞一路推到底，比有空气时推得更干脆。选（b）。

\textbf{热水为什么能顶走活塞？}热水把针筒内空气的温度抬高，\(T\) 上升，分子平均动能变大，在原来的体积下内部压强超过了大气压，碰撞之差把活塞向外推；活塞外移、体积变大，压强回落，直到重新平衡。冷却则反过来。热变成推力的链条是：热 \(\rightarrow\) 分子平均动能 \(\rightarrow\) 碰撞变狠 \(\rightarrow\) 压强差 \(\rightarrow\) 推力。这条链条也是一切热机的命脉，第 28 章会把它做成正经的“能量总账”。

\textbf{薯片袋为什么在高原上鼓起来？}袋子密封，内部 \(N\)、\(T\) 大致不变，内部压强基本还是平原封装时的一个大气压；可高原上的外部大气压只有平原的七八成。内外碰撞之差把袋子撑鼓——不是里面的空气“变多了”或者“变凶了”，而是外面的空气变稀了。

\textbf{对照实验里水为什么压不动？}液体分子之间几乎紧挨着，推开哪怕一点点都要对抗强大的分子间作用力，体积几乎不变；气体分子间距是自身尺寸的十倍量级，中间全是可以让出的空档，压缩只是让分子住得挤一点，反抗只来自碰撞变密，所以“软”得多。同样是堵住孔推活塞，一个压的是间距，一个压的是碰撞频率。

顺手核销那只吸盘：把它按在瓷砖上时，你把吸盘内的空气挤走了大半，内部压强远低于外部大气压，外侧每平方厘米约十牛的碰撞差把吸盘牢牢按在墙上，挂一串钥匙绰绰有余。“吸盘吸墙”这句话方向恰好反了——是大气压在压它。

最后看一个把本章全部概念串起来的工程应用：高压锅。锅盖密封后持续加热，锅内水不断变成水蒸气，\(N\) 增大、\(T\) 升高，体积又被锅体锁死——三个变量里两个往上顶，压强一路爬升，直到顶起限压阀，稳定在约两个大气压。接下来的事属于第 26 章的相变知识：水的沸点随压强升高而升高，锅里的汤水可以在约 120 摄氏度下保持液态。食物熟的快慢对温度极其敏感，这二十度的跃升把炖肉的时间砍掉一大半。而那个小小的限压阀，本质上是一只被压强差顶着的活塞：内部碰撞减去外部碰撞恰好能把它顶起一条缝时，多余的蒸气放走，压强就此锁定。针筒、吸盘、高压锅、恒星，一条原理贯穿到底。

\step{脑洞实验与挑战题}

\begin{thought}[用气体压强顶住一颗恒星]
把本章的画面推到宇宙尺度。太阳是一个直径一百四十万千米的气体球，没有任何固体表面，它的物质在自身引力下本该一路向中心坍缩。是什么顶住了引力，让太阳五十亿年保持这个个头？答案之一是压强：太阳核心的温度高达上千万开尔文，粒子平均动能巨大，向外的碰撞压力层层向外传递，每一层气体都被“下面撞得狠、上面撞得轻”的压强差托住，正好平衡头顶上全部物质的重量。这和针筒里活塞被碰撞之差顶住，是同一条物理。当然太阳还有更深层的故事——核心的核聚变不断补充能量，维持这份温度；等到燃料耗尽、压强顶不住的那一天，恒星会坍缩，演出白矮星、中子星乃至黑洞的后传。这串后传属于本书后面的章节。现在只需要记住：顶住太阳的和顶回你针筒活塞的，是同一种“隐形的推手”。
\end{thought}

\begin{puzzles}
\item \textbf{两只气球。}回到第二题：完全相同的两只气球，充到一样大，一只装空气，一只装氦气，温度相同。哪只内部压强大？提示：气球大小相同意味着体积相同、膜的绷紧程度相近；压强取决于 \(N\)、\(V\)、\(T\)，不取决于分子质量。再想一层：同温度下氦分子比氮分子轻得多，谁的平均速率大？既然撞得更轻，为什么压强能一样？（质量小的分子用更高的速度补回来——\(\frac{1}{2}m\overline{v^2}\) 只由温度决定。）

\item \textbf{夏天的轮胎。}汽车手册建议按“冷胎”胎压打气，并警告不要把气瓶和充气过足的轮胎放在烈日下暴晒。用状态方程解释：暴晒为什么危险？提示：轮胎体积几乎不变、分子数不变，\(p\) 正比于 \(T\)——注意要用开尔文温度估算比例：从 300 开尔文晒到 340 开尔文，压强上涨约百分之十三，而轮胎和瓶盖的安全余量是有限的。

\item \textbf{冰箱里的气球。}把一只扎紧的小气球放进冰箱冷藏室，一小时后观察，再拿出来放一小时。预测并解释两个阶段的体积变化。如果气球里装的是水蒸气含量很高的热空气（比如对着气球呼出的气），冷缩会比纯空气更明显，为什么？提示：分两条线索想——温度变化改变 \(T\)；而水蒸气遇冷凝结成水，直接减少了气相分子数 \(N\)。哪条线索主导，取决于初始湿度。

\item \textbf{针筒定量版。}堵住孔的针筒，初始体积 10 毫升，内外都是 1 个大气压。你缓慢推到 5 毫升并保持（温度不变）。活塞内外侧面的压强差是多少大气压？如果活塞截面积是 2 平方厘米，你的手大约要用多大的力顶住它？提示：压强差约 1 个大气压，即约 \(10^5\) 牛每平方米；力等于压强差乘面积。算出结果后想一想：为什么打气筒打到后面那么费劲——你的对手从来不是橡胶，而是碰撞频率。
\end{puzzles}

本章的图景——用一大群乱撞的粒子解释压强、温度和体积之间的关系——是统计物理的第一次亮相。第 28 章将把它接进能量守恒的总账，讨论压缩和膨胀时气体与外界交换的功和热；第 29 章则要回答一个更不安的问题：如果微观粒子的运动都可逆，为什么打碎的杯子不会自己复原？
